\section{Evaluation}
\label{sec:evaluation}

The central question is whether workflow-aware protection reduces
\emph{workflow completion time} (WCT) after a failure without paying a
larger cost during normal service. We organize the evaluation around
end-to-end outcomes, recovery-path costs, mechanism contributions,
prediction quality, and operating limits. The experimental design follows
three useful precedents: application-level outcomes in workflow-aware
serving~\cite{lin2024parrot,ni2026topas}; separate interrupted and
uninterrupted cohorts in failure recovery~\cite{cao2026lumen}; and
end-to-end measurements followed by restoration, scheduling, and
resource-pressure ablations in KV-cache systems~\cite{gao2025hcache,qin2025mooncake}.

\paragraph{What is measured.}
The current artifact contains physical KV-path microbenchmarks,
held-out predictions on observed agent traces, native sparse-access
captures, and historical analytical diagnostics. It does not yet contain
a fault-injection campaign for the current fixed-$L/K/Z$, exclusive
three-tier controller. We report the available evidence at its actual
scope and specify the missing experiments below. In particular, an older
Peer-disabled allocation model and an isolated rental model are
diagnostics, not implementations of the present policy. The figures
distinguish physical measurements, trace observations, and models; no
synthetic curve stands in for an unexecuted system comparison.

\subsection{Methodology and Comparison Contract}
\label{sec:eval-method}

\paragraph{Deployment and workloads.}
The primary system experiment requires at least two compatible serving
replicas in distinct failure domains, independently reachable backing
storage, and a durable input/output log. GPU-process, replica, host, and
backing-store failures are separate cells. Record DP and TP independently,
the surviving compute capacity, physical interconnect, backend, model
revision, precision, attention kernels, block size, and controller version.
Local NVMe, remote NVMe, and a shared filesystem are distinct backends;
a local disk is not host-failure-safe. These deployment fields remain
unfilled until an executable prototype and hardware allocation are frozen.

Use executable coding, multi-hop question-answering, and controlled
sequential/fork--join workflows. Preserve the original DAG once, complete
prompt histories, workflow IDs, and absolute arrival sequence. Frozen
tool responses provide repeatable primary comparisons; live tools are
a separate validity experiment. Dense and native-sparse model paths are
separate cells, with unsupported context lengths explicitly excluded.
Calibrate one fixed baseline's sustainable healthy load, then apply the
same open-loop arrivals to every policy. The initial decisive slice uses
three workload families, two compatible model paths, one load/budget
setting, four policies, two fault conditions, and ten paired campaigns:
480 planned run slots, not 480 completed runs. Broader load, context,
capacity, and fault-position scans follow the decisive slice.

\paragraph{Baselines and resource accounting.}
Recompute retains accepted tokens and application state but no additional
recoverable KV; it measures the benefit of checkpointing itself.
Hierarchical LRU (H-LRU) and reuse-only protection use the same transport,
checkpoint trigger, recovery scheduler, ownership protocol, and total
physical resources as AegisKV. H-LRU tests whether the new priorities beat
a tuned hierarchy; reuse-only isolates recovery urgency from normal
locality. A zero-slack control isolates workflow urgency while preserving
normal-reuse value. Peer-off and fixed-total-Peer-cap controls isolate
Peer acceleration and the Adapter.
All-DRAM and small-instance offline policies are explicitly labeled
references with their relaxed assumptions. Compatible external
D\'ej\`aVu, GhostServe, LUMEN, or Mooncake artifacts provide additional
system comparisons~\cite{strati2024dejavu,jayakody2026ghostserve,cao2026lumen,qin2025mooncake};
a reimplemented mechanism is named as such, and an incompatible artifact
is not assigned invented performance.

All equal-resource policies pay for execution and protection HBM,
metadata, staging, reserved and retiring bytes, checkpoint writes, and
migration traffic. Shared objects consume one stable physical allocation;
both copies count during handoff. The union of Peer, DRAM, and storage
owners covers a checkpoint. We do not give Peer a free complete lower-tier
replica. $L$, $K$, and $Z$ are selected on calibration workflows and frozen
for held-out evaluation. Adaptive per-request depth or hot-set size is
a separate reference; the current Adapter already varies total Peer
capacity while keeping the per-request quotas fixed.

\paragraph{Outcomes and statistical unit.}
For policy $p$ and workflow $w$, use that policy's own fault-free mate:
\begin{equation}
 \begin{aligned}
 D_{w,p}&=\mathrm{WCT}^{F1}_{w,p}-\mathrm{WCT}^{F0}_{w,p},\\
 G_{w,p}&=D_{w,\mathrm{H\mbox{-}LRU}}-D_{w,p}.
 \end{aligned}
 \label{eq:eval-penalty}
\end{equation}
$F0$ denotes no injected failure; $F1$ denotes the matched fault condition.
Positive $G$ indicates a smaller failure penalty. Report absolute WCT,
completion by a common horizon, SLO goodput, and healthy-service cost as
well as $D$ and $G$. Separate directly interrupted workflows, workflows
sharing recovery resources, and an isolated control cohort. Campaigns
are the independent repetition unit; workflows, requests, and blocks
within a campaign are not independent runs. Report both object counts
and campaign counts, and aggregate paired differences within each
campaign before computing campaign-level uncertainty.
TTFT starts at the original request arrival, including retries. For
requests already decoding at failure, report failure-to-first-resumed-token
latency and the inter-token latency tail. End-to-end TPOT includes recovery
pauses between the first and last accepted output tokens; decoder-only
TPOT is a separately labeled diagnostic. Replayed or discarded tokens do
not enlarge the accepted-output denominator.

\subsection{End-to-End Workflow Completion}
\label{sec:eval-workflow}

\paragraph{Decisive experiment.}
The primary fault is triggered when a fixed workflow/Agent/request reaches
80\% of its frozen logical token progress. Kill the complete source
replica and include every interrupted request. Sweep 20/50/80/95\%
progress separately; do not substitute one wall-clock timestamp for
equal progress. In attribution runs, freeze the pre-fault source mapping
and committed coverage. In full-system runs, allow routing and achieved
frontiers to differ, but compare the original arrival cohort and report
the actual interrupted IDs and wall-clock trigger times. Do not select
a different interrupted subset for each policy. Routing and trigger-time
differences remain part of the full-system effect; mechanism attribution
uses the frozen mapping and coverage control above.

The primary result should pair a WCT distribution with same-ID $G$
values for each workload family and load. Show completion and timeouts
alongside latency; use a shared eligible cohort for complete-case
latency. Failure penalties can be zero or negative, so ratios with $D$
as the denominator are inappropriate. A convincing result must improve
over H-LRU or reuse-only, not only over Recompute, and must survive
accounting for the healthy-service cost.

\paragraph{Available allocation diagnostic.}
Figure~\ref{fig:eval-placement-pairs} reanalyzes the existing
36-configuration, four-branch model. It varies state size, branch slack,
and storage bandwidth while comparing identical configurations. The
historical workflow-value allocator improves 17 cases and ties 19. In
the base case, LRU and workflow value add 3.6208 and 1.6125\,s to a
12\,s residual join, respectively. The 2.0083\,s difference is a
55.5\% reduction in the \emph{failure penalty}, not in full-workflow WCT.
This model has no Peer tier, fixed-slot admission, or lazy lower-tier
controller. It supports the relevance of workflow priorities and the
choice of a paired metric, but not an end-to-end AegisKV speedup.

\begin{figure}[tb]
 \centering
 \includegraphics[width=\columnwidth]{figures/evaluation/eval-placement-pairs.pdf}
 \caption{Historical allocation model: paired reduction in failure penalty
 relative to LRU across 36 equally weighted configurations. Positive values
 favor workflow-value placement; 19 ties remain visible. These are
 deterministic Peer-disabled scenarios, not measured workflows.}
 \label{fig:eval-placement-pairs}
 \Description{Paired LRU minus workflow-value failure penalties over all
 36 historical analytical configurations, preserving positive gains and ties.}
\end{figure}

\subsection{Recovery Paths and Interference}

\paragraph{Available physical calibration.}
Figure~\ref{fig:eval-recovery-paths} reorganizes the existing path timings
used in the earlier reconstruction comparison. The configuration is Qwen2.5-7B-Instruct-1M,
two A100 PCIe 40\,GB GPUs with TP2, BF16, and batch size one. One fixed
synthetic token stream is evaluated at 8K, 64K, 256K, and approximately
1.01M tokens. HBM/DRAM reuse use three timing repetitions per input;
full Prefill and DPC shared-filesystem reuse use one. The endpoint is
the same input KV ready for continuation, including retrieval and
required computation. These are path measurements, not independent
workflow campaigns.

At the longest context, resident HBM, host DRAM, shared filesystem,
and full reconstruction take 1.3050, 2.6632, 10.1561, and 571.2127\,s.
The large reconstruction gap motivates retained state; the smaller
HBM--DRAM gap is the quantity a Peer tier could potentially reduce.
This measurement does not include a Peer transfer, background serving,
concurrent recovery, or fault detection. The filesystem server cache is
uncontrolled, so this is neither cold-SSD latency nor a bandwidth claim.

\begin{figure}[tb]
 \centering
 \includegraphics[width=\columnwidth]{figures/evaluation/eval-recovery-paths.pdf}
 \caption{Measured KV-path latency on one fixed input stream. HBM and
 DRAM show medians and observed ranges of three timing repetitions;
 reconstruction and shared-filesystem points each have one observation.
 HBM is resident reuse, and SharedFS is DPC shared storage. These are
 calibration measurements, not fault-injected system results.}
 \label{fig:eval-recovery-paths}
 \Description{Four measured recovery-related path latencies across four
 context lengths, with observed timing ranges only where repetitions exist.}
\end{figure}

\paragraph{Required causal decomposition.}
Log detection, fencing/rerouting, admission, directory lookup, storage
read, network transfer, H2D, required replay, request resume, and DAG
completion. Trace dependency edges and identify the realized critical
path; adding overlapping component durations would overcount recovery.
Hold demand and coverage fixed to compare Peer off/on with overlap
off/on. Separately measure total bytes, unused prefetch, GPU work,
temporary copies, and recovery-pool busy time. A faster request does not
necessarily reduce join delay or total recovery work.

Inject bursts affecting 1/2/4/8 workflows with 1/2/4 surviving targets
under the same foreground arrivals. Plot queue occupancy and
interrupted/uninterrupted WCT, TTFT, and TPOT over time, with an isolated
resource control. Shared recovery queues must be identified before
claiming collateral delay. This follows the cohort separation in
LUMEN~\cite{cao2026lumen} while retaining WCT as AegisKV's primary outcome.

\subsection{Mechanism Contributions and Credit Cost}
\label{sec:eval-rental}

Use matched interventions rather than removing several mechanisms at
once. The main ablations are: zero-slack fallback versus predicted slack;
reuse-only versus recovery-only versus combined priority; per-Agent
shares versus creator-only shares with physical deduplication unchanged;
causal-visibility normalization versus
raw selection counts; fixed-slot Slack versus FCFS, recency, and
recovery-cost admission; and fixed Peer capacity versus guarded AIMD.
For lower tiers, compare the current lazy-credit policy with same-value
direct revaluation, ordinary cost/size GDS or Landlord, no actual-use
renewal, and no pressure aging. Preserve the same ownership, leases,
arrival/fault opportunities, value estimates, and physical budgets.

Lazy versus eager refresh needs both an equal-CPU-budget comparison and
an unbounded eager reference. Report decision time, candidates examined,
scan age, stale rankings, differing victims, movement bytes, and paired
WCT. Fewer scalar checks alone do not establish a faster controller.
Five factorial experiments expose interactions: Slack $\times$ shared-Agent
shares, normalized demand $\times$ restoration pipelining, renewable credit
$\times$ lazy refresh, load-aware routing $\times$ shared-link contention,
and Peer residency $\times$ pipelining. Each uses all four cells with
the other features fixed. Campaign-level differences of differences
distinguish a combined effect from two separately observed improvements.

\paragraph{Available adverse result.}
The historical additive rental diagnostic has 512 decisions and the
same 396 value checks for direct revaluation and the original trial-rent
variant. Their miss-plus-movement costs are 462.96 and 506.39 model units:
trial rent is 9.38\% worse. Figure~\ref{fig:eval-rental-diagnostic} also
retains the later repaired variant. A diagnostic repair does not validate
the current per-Agent lazy-credit controller. This counterexample makes
the same-value direct baseline necessary; it is not evidence that adding
rent improves WCT or reduces CPU cost.

\begin{figure}[tb]
 \centering
 \includegraphics[width=\columnwidth]{figures/evaluation/eval-rental-diagnostic.pdf}
 \caption{Historical rental diagnostic. Miss and movement costs are
 retained for direct revaluation, the original trial-rent variant, and
 its repaired variant. Lower is better. The original negative result is
 preserved; model cost and value-check counts are not physical WCT or
 controller CPU measurements.}
 \label{fig:eval-rental-diagnostic}
 \Description{Historical miss-plus-movement cost comparison, preserving
 the rental regression and the remaining gap after a diagnostic repair.}
\end{figure}

\subsection{Prediction Quality and Sparse Demand}
\label{sec:eval-prediction}

\paragraph{Remaining time.}
Figure~\ref{fig:eval-tracker-errors} uses 1,443 eligible second-request
agent snapshots from 170 workflows, with one snapshot per Agent and
five workflow-disjoint folds. The target is whole-Agent remaining wall
time, not the residual segment duration needed by the current Tracker.
Mean absolute errors are 438.6\,s for a training-fold median, 406.7\,s
for basic GBDT, and 407.1\,s for GBDT with history. Basic GBDT's p95
absolute error is nevertheless worse than the median baseline
(1113.8 versus 1036.1\,s). The plot preserves
the complete error tail, including errors above 85,000\,s. A modest
average reduction does not establish trustworthy slack, and these
labels cannot validate segment-level predictions.

The current Tracker requires new causal segment labels, workflow-disjoint
training/calibration/test sets, and later-workload validation. Measure
remaining-time error, false positive slack on truly critical branches,
critical-path ranking, fallback frequency, and predictor/update cost.
Evaluate zero-slack and oracle-duration references. Bootstrap whole
workflows for prediction diagnostics; do not infer an error CDF from
aggregate WAPE or feature correlations. Study stale profiles, model
shifts, missing DAG edges, and delayed tools by recomputing decisions
against actual outcomes, rather than adding decorative error bars.

\begin{figure}[tb]
 \centering
 \includegraphics[width=\columnwidth]{figures/evaluation/eval-tracker-errors.pdf}
 \caption{Held-out absolute error of whole-Agent remaining-time
 predictions. The ECDF contains all 1,443 Agent snapshots from 170
 workflows; folds separate workflows. It is descriptive and has no
 invented confidence band. Whole-Agent accuracy does not validate
 residual-segment slack.}
 \label{fig:eval-tracker-errors}
 \Description{Complete held-out absolute-error ECDFs for median, basic
 GBDT, and history-augmented GBDT, including the long error tail.}
\end{figure}

\paragraph{Native access and horizon sensitivity.}
Figure~\ref{fig:eval-sparse-demand} compares historical frequency,
recency, and an offline demand oracle at the same within-model byte
budget. The captures use 32,768 historical tokens, a preceding
512-Query ranking window, and the following 64 Queries for evaluation.
Each object is a layer's 64-token block including all KV heads. The
budget is 205 of 512 historical blocks per layer, approximately 40\%.
Coverage is the fraction of distinct future-needed historical bytes
already selected, with equal document weighting and three documents
per model. It is neither attention mass nor a GPU hit-rate measurement.

Historical frequency covers 77.82\%, 53.57\%, and 43.02\% for
NOSA-3B, MiniCPM4-8B, and the experimental SeerAttention/Qwen3 path;
same-budget oracle coverage is 93.75\%, 58.98\%, and 43.28\%.
The small Seer oracle gap limits the value of a better ranker at this
budget. Extending the future window to 512 Queries lowers frequency
coverage to 58.90\%, 44.08\%, and 40.47\%; the complete horizon sweep
is retained in the artifact. These historical counts include native
head/group selection events. They are a baseline for, not a measurement
of, the current deduplicated, visibility-normalized probability estimate.

The required current-policy experiment compares raw counts, normalized
counts, recency, and random selection using strictly preceding windows.
An equal-budget offline oracle may use actual future demand and is
reported only as a clairvoyant reference. Sweep history length, future horizon,
budget, and unseen documents/workflows; account for selector state and
all physical objects. Freeze $L/K/Z$ on calibration workloads and test
them on held-out workloads. Finally measure exposed wait and WCT under
misses and foreground contention: high coverage alone does not prove
net Peer benefit.

\begin{figure}[tb]
 \centering
 \includegraphics[width=\columnwidth]{figures/evaluation/eval-sparse-demand.pdf}
 \caption{Native sparse-access diagnostic at an equal 205-of-512-block
 budget per layer. Historical frequency and recency select for the
 following 64 Queries; the offline oracle sees their realized demand.
 Points and summaries retain the
 three-document variation; this is observed demand coverage, not a
 confidence interval or failure-recovery speedup.}
 \label{fig:eval-sparse-demand}
 \Description{Same-budget native demand coverage for three models and
 three placement references, with the individual document variation retained.}
\end{figure}

\subsection{Healthy Service and Sustained Failures}

\paragraph{Normal-service cost.}
Run the complete original workloads under $F0$, including first calls,
tool waits, and follow-up requests. Measure WCT, TTFT, TPOT, SLO
goodput, GPU-seconds, controller CPU/memory, checkpoint bytes, and
per-tier request and token hits. Hold absolute arrival traces fixed at
30/60/85\% of the reference capacity; do not normalize each policy to
its own throughput. Record execution-HBM loss and migration contention
when Peer grows. Prefix overlap in a trace is not a runtime cache hit.
Report conditional recovery gains separately from any amortized benefit
at a declared failure frequency. A model queueing delay is not measured
TTFT, and a failure-heavy stress test does not establish production net
benefit.

\paragraph{Why averages are insufficient.}
The historical LANL-timing model~\cite{lanlFailureData,schroeder2006failures}
retains 21,738 cleaned node-outage
intervals from 22 systems. Each start triggers an assumed four-branch
recovery in a per-system FIFO spare pool. Figure~\ref{fig:eval-trace-tails}
shows a lower event-weighted mean for workflow-value allocation
(9.4673 versus 10.4533\,s) but worse p95 (40.0292 versus 37.0375\,s)
and p99 (232.1125 versus 204.1208\,s). Its longer pool occupancy can
hurt subsequent recoveries even when the current join finishes sooner.
These are modeled outcomes driven by recorded timing, not observed
LANL agent workflows or measured AegisKV failures.

\begin{figure}[tb]
 \centering
 \includegraphics[width=\columnwidth]{figures/evaluation/eval-trace-tails.pdf}
 \caption{Historical trace-driven recovery model: mean and tail failure
 penalties under the same recorded event timing. Workflow-value
 placement reduces the mean while worsening p95 and p99. The outcome
 combines an assumed workload and spare-pool model with recorded outage
 starts; it is not measured fleet availability.}
 \label{fig:eval-trace-tails}
 \Description{Mean and tail recovery penalties in a historical FIFO
 model, retaining the adverse p95 and p99 comparison.}
\end{figure}

For continuous replay, preserve failure clusters, downtime, surviving
capacity, and non-overlapping calendar blocks; report event-weighted
and system/window-weighted summaries. The available LANL redistribution
has 23,739 rows versus 23,741 documented by its original source, an
unresolved provenance difference. Reflex/ByteDance
alerts~\cite{ding2026reflex,reflexDataset} are shifted
within each job and lack reliable down/up pairs, so they cannot define
a global failure calendar. Modern node-fault traces such as
InfiniteHBD~\cite{shou2025infinitehbd} also require their original
node-level semantics and training-cluster scope. Trace provenance and
explicit resource mapping precede any claim about serving availability.

\subsection{Operating Range and Controller Scaling}

Capacity experiments ask when protection loses its advantage. Independent
Peer, DRAM, and storage sweeps include Peer-off, full destinations, and
last-copy pressure under equal physical budgets. Achieved coverage and
commit lag distinguish insufficient capacity from an overloaded copying
path: the uncommitted tail $N$ may exceed the checkpoint interval $Z$.
WCT is therefore interpreted alongside write backlog, migration traffic,
unused prefetch, and fallback frequency. Separate storage, NIC, and H2D
bottlenecks reveal whether a gain depends on one favorable transport ratio.

Scaling experiments separate replica count, TP degree, and failure domains
from controller state size. Increasing DAG size, shared-owner fan-out,
live objects, and candidate budget exposes the tradeoff between decision
quality and update cost. Actual decision-time tails, CPU use, metadata,
and foreground stalls compare bounded lazy refresh with equal-budget
and unbounded eager controls. Saturation and unsupported cells remain
visible; multiplying operation counts by an assumed node cost cannot
substitute for controller timing.

Robustness experiments test whether fallback and capacity control remain
useful when duration errors correlate across a DAG, profiles become stale,
or tools and arrivals change abruptly. Step, ramp, and burst loads expose
Adapter overshoot and settling time. The key capacity outcome is actually
freed bytes: lowering a quota target does not release memory until the
Peer-to-DRAM-to-storage cascade and reader drain finish.

\subsection{Recovery Correctness and Reproducibility}
\label{sec:eval-plan}

Correctness is a prerequisite for interpreting any latency result. The
state-transition suite checks exclusive coverage, commit-before-redirect,
and reader-drain-before-free under stale epochs, duplicate events, and
incomplete restore state. Sole-owner loss must lead to a declared
older-checkpoint/replay fallback. Pressure and tier-round-trip cases
check physical occupancy, retained paid rent, and cancellation; Peer
cases check the specified victim order, tie rules, and physical leases.
The artifact enumerates the injection points and seeded interleavings.

Faults bracketing copy, commit, ownership switch, and acknowledgement
test accepted-token durability and complete compatible restore state.
Next-output equivalence uses a declared numerical tolerance, while
acknowledged output must never be duplicated. Application tool
idempotency is a separate requirement. Transition coverage and failed
randomized schedules accompany the results; earlier replicated-Peer
examples cannot establish these invariants for exclusive ownership.

Freeze source/model/backend versions, workload and fault IDs, physical
budgets, and calibration choices before collecting primary results.
Randomize policy order and start with ten paired campaigns; increase
repetition based on pilot variance and precision rather than the desired
sign of an effect. Aggregate by campaign and publish bootstrap method,
seed, and intervals. Deterministic model scenarios have no experimental
confidence intervals; the microbenchmark ranges describe repeated
timings of one input, not independent workflow uncertainty. Preserve
timeouts in the completion denominator and do not invent finish times.

The artifact contains the complete experiment matrix, per-figure source
hashes, long-form plotting tables, native-size vector exports, and
generation commands. All unexecuted system cells remain explicitly
planned. The execution order is correctness and backend calibration,
paired single-failure outcomes, healthy overhead and decisive ablations,
then continuous trace replay and scale/pressure tests. Attractive plots
and a complete protocol do not close the current gap in end-to-end
implementation evidence.
